\section{Replicating a Vacuum Robot Simulation}

This is a worksheet for the first SF35803  lab.


Imagine you are a software developer working under my supervision. I have recently come across an exciting demonstration of a vacuum robot in action which is accessable via the \href{https://youtube.com/shorts/G1tke2bdzCo}{link}.  Impressed by its functionality, I want you to replicate a simplified simulation of this robot's behavior.

\subsection{Scenario 1: Initial Attempt - Can LLMs Do It All?}


As a first step, I want to see how effectively you can utilize readily available Large Language Models to achieve this task, even if you have limited to no prior coding experience.
\subsubsection{Your Task: Explore LLM Capabilities with a Simple Prompt}



\begin{enumerate}
    \item \textbf{Choose an LLM:} Select any Large Language Model platform you are comfortable with, such as ChatGPT, Gemini, DeepSeek Coder, or any other similar service.

    \item \textbf{Observe the Vacuum Robot:} Carefully watch the \href{https://youtube.com/shorts/G1tke2bdzCo}{provided video}
    \item \textbf{Formulate an Initial Prompt:} Based on your observation of the vacuum robot's behavior, write a prompt for your chosen LLM. This prompt should be as natural and straightforward as possible, reflecting an initial attempt by someone with limited coding knowledge who believes in the "magic" of AI.  For example, you might write something like:

    \begin{sloppypar}
    {\ttfamily\raggedright Write code to simulate a vacuum robot moving in a room and cleaning dirt.\par}
    \end{sloppypar}



    \item \textbf{Generate Code:} Input your prompt into the LLM and request it to generate the code. Since I am comfortable in Python, you should restrict to provide the output code in Python. But, you may ask how to run the code tho and ask the LLM to also suggest any necessary libraries it deems appropriate for this task.

    \item \textbf{Execute the Code:} Attempt to execute the code generated by the LLM. Note down any errors, issues, or the level of success you achieve.

    \item \textbf{Analyze the Outcome:} Reflect on the results. Did the LLM produce functional code? How closely does it resemble a vacuum robot simulation? What are the shortcomings of this initial approach?

\end{enumerate}

% \subsection*{Expected Outcome of Scenario 1:}

% It is expected that your initial, simple prompt might not yield a fully functional or detailed simulation. This scenario is designed to simulate a situation where relying solely on a basic prompt might not be sufficient to achieve complex software development goals, even with powerful LLMs.

\newpage
\subsection{Scenario 2: Detailed Instructions - A Helping Hand}

Let's assume your initial attempt in Scenario 1 did not produce the desired, detailed vacuum robot simulation. As your supervisor, recognizing the potential of LLMs but also the need for more specific guidance, I am now providing you with a much more detailed instruction set in the form of a comprehensive prompt.

\subsubsection{Your Task: Refine Simulation with a Detailed Prompt}


So to begin,
\begin{enumerate}
    \item \textbf{Access Detailed Prompt:} Navigate to the following link to access a detailed prompt specifically designed for generating a vacuum robot simulation: \href{https://github.com/balandongiv/computer_programming_gist/blob/630ed5a5dfb33c305fb60840889597d7f8b0dd94/session1/robot_simulation_prompt.md}{Detailed Prompt Gist}.

    \item \textbf{Understand the Detailed Prompt:} Carefully read and understand the detailed prompt provided. Notice how it breaks down the simulation into components, specifies behaviors, and describes visual elements.

    \item \textbf{Use the Detailed Prompt with LLM:}  Paste the entire detailed prompt into the same LLM you used in Scenario 1 (or you may choose to try a different LLM). Request it to generate code based on this detailed prompt.

    \item \textbf{Execute the Code (Again):} Attempt to execute the code generated using the detailed prompt.  Pay close attention to any instructions provided by the LLM regarding setup, libraries, or execution.

    \item \textbf{Finger Cross}: Do you manage to replicate the robot movement?


\end{enumerate}

\newpage
\subsection{Reflection: LLMs as Tools, Foundational Skills as the Foundation}

This tutorial aimed to provide you with a hands-on experience exploring the capabilities and limitations of Large Language Models (LLMs) in a practical coding scenario.  As you've likely discovered through Scenarios 1 and 2, while LLMs are undeniably powerful and can generate code snippets with remarkable ease,  relying on them \textit{solely}, especially without a foundational understanding of programming, presents significant challenges.

Let's reflect on what you might have encountered, particularly if you approached this as someone new to coding:

\subsubsection{Scenario 1: The "Magic" Isn't Always Enough}
Your initial, simple prompt likely resulted in code that, at best, was a very basic interpretation of a vacuum robot's movement.  It may have lacked crucial features, exhibited unrealistic behavior, or even been non-functional. This illustrates that LLMs aren't mind-readers.  They require clear, specific instructions to generate code that aligns with complex requirements.  The "magic" is in their ability to process language, but the quality of the output is heavily dependent on the quality of the input and the underlying understanding of the task.

\subsubsection{Scenario 2: Detail Improves Output, But Doesn't Solve Everything}
The detailed prompt in Scenario 2 undoubtedly led to a more sophisticated and potentially functional simulation.  You likely saw improvements in the robot's behavior and perhaps even visual elements.  However, even with a detailed prompt, you might have still faced hurdles.  Perhaps the generated code:
\begin{itemize}
    \item \textbf{Contained errors:} LLMs can still produce code with syntax errors or logical flaws, even with detailed instructions.
    \item \textbf{Wasn't fully complete:}  The LLM might have missed certain aspects of the desired behavior, or made assumptions that weren't explicitly stated in the prompt.
    \item \textbf{Was difficult to understand and modify:} Even if the code \textit{ran}, you might have found it challenging to decipher \textit{how} it worked, making it difficult to customize or debug.
    \item \textbf{Required external libraries you weren't familiar with:}  LLMs might suggest using libraries you've never encountered, adding another layer of complexity to getting the code to run.
    \item \textbf{Presented execution challenges:} Setting up the correct environment, installing libraries, and running Python code in general can be daunting for beginners.
\end{itemize}


\newpage
\subsection{Common Newbie Experiences Relying Solely on LLMs}

If you were approaching this without prior coding experience, you might have encountered frustrations such as:

\begin{itemize}
    \item \texttt{"It doesn't work!"} without understanding \textit{why}:  Errors might appear without clear explanations, leaving you unsure how to troubleshoot.
    \item \textbf{Feeling lost in the code:} Even if the code runs, you might not grasp the underlying logic or how different parts interact.
    \item \textbf{Struggling to install libraries or set up the environment:} The practical aspects of running code are often overlooked by simple prompts.
    \item \textbf{Being unable to adapt or extend the code:} If you wanted to add a new feature or fix a bug, you might feel helpless without understanding the code's structure.
    \item \textbf{Over-reliance and eventual disappointment:}  Initially impressed by the LLM's output, you might become disillusioned when faced with the practical challenges of making it work and customizing it.
\end{itemize}

\newpage
\subsection{The Crucial Role of Foundational Programming Skills}

This experience highlights that while LLMs are powerful \textit{tools} for code generation and can be incredibly helpful for experienced programmers, they are \textit{not} a replacement for fundamental programming knowledge.  Think of an LLM like a sophisticated power tool – it can be incredibly efficient in the hands of a skilled crafts person, but potentially dangerous and ineffective for someone without basic carpentry skills.

This course is designed to equip you with those essential "carpentry" skills for programming.  Through the topics we will cover, you will gain:

\begin{itemize}
    \item \textbf{Systematic Thinking and Algorithmic Formulation:}  Learn to break down problems into logical steps, enabling you to create effective prompts and understand the underlying logic of code, whether generated by you or an LLM.
    \item \textbf{Python Fundamentals:}  Master the basic building blocks of programming in Python, allowing you to read, understand, and write code independently.
    \item \textbf{Control Flow (Repetition and Decision Structures):}  Learn how to create dynamic and interactive programs, going beyond simple linear code generation.
    \item \textbf{Functional and Object-Oriented Programming:}  Develop modular and reusable code, making your projects more manageable and scalable, and enabling you to better understand more complex LLM-generated code.
    \item \textbf{Data Structures (Arrays) and Data Visualization:} Learn to work with data effectively and present it visually, expanding the scope of projects you can undertake.
    \item \textbf{Error Management and Debugging:}  Develop essential skills to identify, diagnose, and fix problems in your code, and in code generated by LLMs, making you a more resilient and independent programmer.
    \item \textbf{Test-Driven Development:} Learn to write robust and reliable code by incorporating testing methodologies.
\end{itemize}

\newpage
\subsection{Course Learning Outcomes Revisited}

By mastering these topics, you will directly achieve the course learning outcomes:

\begin{itemize}
    \item \textbf{Apply the process, methods and fundamental knowledge in designing algorithms and the concepts of programming logic:} You'll move beyond simply prompting LLMs and gain the ability to design algorithms yourself, understanding the "why" behind the code.
    \item \textbf{Construct the logic and algorithm design to represent logic in high-level programming language:} You'll be able to translate your algorithmic thinking into actual Python code, whether writing it yourself or effectively guiding an LLM.
    \item \textbf{Demonstrate the ability to gain new knowledge of related entrepreneurship skills and share with others:}  A solid programming foundation empowers you to explore entrepreneurial opportunities in technology and effectively communicate your knowledge to others.
\end{itemize}

\vspace{40pt}

\begin{quote}
    ``LLMs are impressive, but they're not a magic bullet. They're tools, and like any tool, they have their limitations.''

    \hfill \href{https://www.reddit.com/r/learnprogramming/comments/1f3vo44/llm_are_just_tools/?utm_source=share&utm_medium=web3x&utm_name=web3xcss&utm_term=1&utm_content=share_button}{LogicalBum}
\end{quote}